\begin{proof}
%Let $\boldsymbol{B}=boldsymbol{B}(k,t,c)$. 
We first calculate 
the characteristic polynomial of $\boldsymbol{B}(k,t,c)$. 
The polynomials $F_i$, $G_i$, $\mathscr{F}_i$, and $\mathscr{G}_i$ are defined in Section~\ref{sec:pre}. 
Note that $F_i(x)$ is the characteristic polynomial of the principal $i \times i$ matrix formed by the first 
$i$ rows and $i$ columns of $\boldsymbol{B}(k,t,1)$ for $t>i+1$. 
By this fact and equations~\eqref{eq:3-term2} and~\eqref{eq:G_i}, we can compute 
\begin{multline*}
|x \boldsymbol{I}-\boldsymbol{B}(k,t,c)|\\
\begin{aligned}
&={%\small 
\begin{vmatrix}
x & -k & & & & \\
-1 & x & -(k-1) & & & \\
 & \ddots & \ddots & \ddots & & \\
 & &-1 & x & -(k-1)& \\
 & & & -c & x & -(k-c) \\
 & & & & -k & x 
\end{vmatrix}}\\
&=k {%\small 
\begin{vmatrix}
x & -k & & & \\
-1 & x & -(k-1) & & \\
 & \ddots & \ddots & \ddots & \\
 & &-1 & x & 0 \\
 & & & -c & -(k-c) \\ 
\end{vmatrix}}+x{%\small 
\begin{vmatrix}
x & -k & & & \\
-1 & x & -(k-1) & & \\
 & \ddots & \ddots & \ddots & \\
 & &-1 & x & -(k-1) \\
 & & & -c & x 
\end{vmatrix}}\\
&=-k(k-c)F_{t-2}(x)+x(x F_{t-2}(x)-c(k-1)F_{t-3}(x))\\
&= c\big(F_{t-2}(x)-(k-1)^2 F_{t-4}(x)\big)+(x^2-k^2)F_{t-2}(x)\\
&= c\big((x^2-k^2)F_{t-4}(x)+F_{t-4}(x)-(k-1)^2 F_{t-6}(x)\big)+(x^2-k^2)F_{t-2}(x)\\
&= (x^2-k^2)\left(c\sum_{i=0}^{ \lfloor (t-4)/2 \rfloor}F_{t-4-2i}(x)+F_{t-2}(x)\right).
%\\
%&=(x^2-k^2)\big((c-1){G}_{t-4}(x)+{G}_{t-2}(x)\big)
\end{aligned}
\end{multline*}
Note that
\[
c\sum_{i=0}^{ \lfloor (t-4)/2 \rfloor}F_{t-4-2i}(x)+F_{t-2}(x)=(c-1){G}_{t-4}(x)+{G}_{t-2}(x). 
\]
Since the zeros of $\mathscr{G}_{\epsilon,s-1}$ and $\mathscr{G}_{\epsilon,s}$ interlace on $(0,4(k-1))$, each zero of $(c-1)\mathscr{G}_{\epsilon,s-1}+\mathscr{G}_{\epsilon,s}$ is simple and belongs to $(0,4(k-1))$ except for the smallest zero. 
For $c=k$ the smallest zero is equal to $0$ because 
$
(k-1)\mathscr{G}_{\epsilon,s-1}(0)+\mathscr{G}_{\epsilon,s}(0)=0
$ 
by~\eqref{eq:new_3-term} and~\eqref{eq:def_scrG}. 
For $c>k$, the smallest zero is negative. From $x^{\epsilon}(x^2-k^2)((c-1)\mathscr{G}_{\epsilon,s-1}(x^2)+\mathscr{G}_{\epsilon,s}(x^2))=(x^2-k^2)((c-1){G}_{2s-2+\epsilon}(x)+{G}_{2s+\epsilon}(x))$, 
each non-zero real eigenvalue of 
$\boldsymbol{B}(k,t,c)$ has multiplicity 1, 
and if $c>k$, then $\boldsymbol{B}(k,t,c)$ has 
imaginary eigenvalues. 

Let $\mathfrak{f}_1(x)$ be the polynomial 
\[
\mathfrak{f}_1(x)=\frac{((c-1){G}_{t-4}(x)+G_{t-2}(x))^2}{x^2-\theta^2}=\sum_{i=0}^{ t- 3} f_i \mathscr{F}_{0,i}(x^2).
\]
We show that $\mathfrak{f}_2(x)=\sum_{i=0}^{ t - 3} f_i \mathscr{F}_{0,i}(x)$ satisfies the condition of the linear programming bound from Theorem~\ref{thm:lp_bound} for bipartite graphs. 
Note that $\mathfrak{f}_2(k^2)=\mathfrak{f}_1(k)>0$, and $\mathfrak{f}_2(\lambda^2)=\mathfrak{f}_1(\lambda) \leq 0$ for each $\lambda \in [-\theta, \theta]$. 
It suffices to show that $f_i>0$ for each $i\in \{0,1,\ldots, t- 3\}$. 
 
The polynomial $\mathfrak{f}_1(x)$ can be expressed by 
\begin{align*}
\mathfrak{f}_1(x)&\Mk =\Mk 
\frac{(c-1){G}_{t-4}(x)+{G}_{t-2}(x)}{x^2-\theta^2}
\left(c\sum_{i=0}^{\lfloor t/2 \rfloor-2}{F}_{t-4-2i}(x)+{F}_{t-2}(x) \right)\\
&\Mk =\Mk x^{2\epsilon}
\frac{(c\Mk -\Mk 1)\mathscr{G}_{\epsilon,\lfloor t/2 \rfloor-2 }(x^2)\Mk +\Mk \mathscr{G}_{\epsilon,\lfloor t/2 \rfloor-1 }(x^2)}{x^2-\theta^2}
\Mk \left(\Mk c\sum_{i=0}^{\lfloor t/2 \rfloor-2}\Mk \mathscr{F}_{\epsilon,i}(x^2)\Mk +\Mk 
\mathscr{F}_{\epsilon,\lfloor t/2 \rfloor-1}(x^2)\Mk \right)\Mk , 
\end{align*}
where $\epsilon=0$ if $t$ is even, 
and $\epsilon=1$ if $t$ is odd. Thus, 
\[
\mathfrak{f}_2(x)=x^{\epsilon}
\frac{(c-1)\mathscr{G}_{\epsilon,\lfloor t/2 \rfloor-2 }(x)+\mathscr{G}_{\epsilon,\lfloor t/2 \rfloor-1 }(x)}{x-\theta^2}
\left(c\sum_{i=0}^{\lfloor t/2 \rfloor-2}\mathscr{F}_{\epsilon,i}(x)+
\mathscr{F}_{\epsilon,\lfloor t/2 \rfloor-1}(x)\right).
\]

By Proposition~3.2 in~\cite{CK07}, 
$g(x)=((c-1)\mathscr{G}_{\epsilon,\lfloor t/2 \rfloor-2}+\mathscr{G}_{\epsilon,\lfloor t/2 \rfloor-1})/(x-\theta^2)$ has positive coefficients 
in terms of $\mathscr{G}_{\epsilon,0},\mathscr{G}_{\epsilon,1},\ldots, \mathscr{G}_{\epsilon,\lfloor t/2 \rfloor-2}$. This implies that $g(x)$ has positive 
coefficients in terms of $\mathscr{F}_{\epsilon,0},\mathscr{F}_{\epsilon,1},\ldots, \mathscr{F}_{\epsilon,\lfloor t/2 \rfloor-2}$.
Therefore $f_i>0$ for each $i=\{0,1,\ldots, t-3 \}$ by 
Lemma~\ref{lem:coe}. 

The polynomial $g(x)$ can be expressed by $g(x)\Mk =\Mk \sum_{i=0}^{\lfloor t/2 \rfloor-2}g_i \mathscr{F}_{\epsilon,i}(x)$. 
By Lemma~\ref{lem:coe}, we have
\[
f_0=\sum_{i=0}^{\lfloor t/2 \rfloor-2} c k^{\epsilon} g_i \mathscr{F}_{\epsilon,i}(k^2)=c k^{\epsilon} g(k^2).
\]
By applying Theorem~\ref{thm:lp_bound} to the polynomial $\mathfrak{f}_2(x)$, we have
\begin{align*}
b(k,\theta) &\leq \frac{2\mathfrak{f}_2(k^2)}{f_0}=2k^\epsilon\left( \sum_{i=0}^{\lfloor t/2 \rfloor-2}\mathscr{F}_{\epsilon,i}(k^2)+\mathscr{F}_{\epsilon,\lfloor t/2 \rfloor-1}(k^2)/c \right) \\
&=2\left( \sum_{i=0}^{t-4} (k-1)^{i}+ \frac{(k-1)^{t-3}}{c}+
\frac{(k-1)^{t-2}}{c}
\right).
\end{align*}
